\documentclass[11pt]{article}
\usepackage{amsmath,amssymb,amsthm}
\usepackage[margin=1in]{geometry}
\newtheorem{theorem}{Theorem}
\newtheorem{lemma}[theorem]{Lemma}
\newtheorem{corollary}[theorem]{Corollary}
\newcommand{\C}{\mathcal C}
\newcommand{\rad}{\operatorname{rad}}
\title{The contraction form of a normalized Cayley cross product}
\author{Research notes for the F004 proof programme}
\date{}
\begin{document}
\maketitle

This note proves the extra nondegeneracy hypothesis needed to apply the
reviewed involution obstruction to a normalized Cayley cross product. It
uses a symbolic Cayley frame, not the coefficients of any retained Original
model. Its final application to a saved three-form still requires the exact
constructor or equality identifying that form with the normalized form below.

\section{Conventions and the classical frame input}

Let $k$ have characteristic different from $2$, and let $\C$ be a
nondegenerate unital Cayley composition algebra over $k$, with multiplicative
quadratic norm $N$ normalized by $N(1)=1$. In this note
\[
 B(x,y)=\frac{N(x+y)-N(x)-N(y)}{2}
\]
is the \emph{half}-polar form. Thus $B(x,x)=N(x)$. This differs by a factor
of two from the full polar form denoted $B$ in the retained Cayley-frame
notes. It is the form denoted $B=\kappa\beta$ in the retained normalization
criterion. Put $V=1^\perp$, so $\dim V=7$, and write the product of pure
elements as
\[
 xy=-B(x,y)1+T(x,y),\qquad T(x,y)\in V.
\]
The normalized cross-product three-form is
\[
 \Phi(x,y,z)=B(T(x,y),z).
\]
It is alternating. This is also immediate from the frame multiplication
below; in the retained normalization construction it is the given alternating
form $\kappa\tau$.

We use the classical Cayley-frame fact already retained in
\emph{Eight rational classes of ordered alternating frames in an
odd-characteristic Cayley algebra}, Section1, exact TeX edition with
SHA256 \texttt{89aaf643b59bb14cacc0d5a391ff48c6b4c51cdc9775d15d819db260db7c9207}.
It is attributed there to the retained Wilson/eigenframe inputs, not proved
as a new recognition theorem: orthogonal nonisotropic pure $i,j$ generate a
nondegenerate quaternion algebra $H=\langle1,i,j,ij\rangle$, and a
nonisotropic $\ell\in H^\perp$ gives the Cayley basis
\[
 1,\ i,\ j,\ ij,\ \ell,\ i\ell,\ j\ell,\ (ij)\ell.
\]
The multiplication is determined by the three norms and this bracketing.
The explicit convention for the norm-one case used here is
\begin{equation}\label{eq:double}
 (a+b\ell)(c+d\ell)
   =(ac-\bar d b)+(da+b\bar c)\ell\qquad(a,b,c,d\in H).
\end{equation}
Thus $i^2=j^2=\ell^2=-1$, $ij=-ji$, and the norm is
$N(a+b\ell)=N(a)+N(b)$. The bar is quaternion conjugation.
This is a choice of standard Cayley-frame convention; no claim that an
Original matrix basis already has this convention is made.

After extending to an algebraic closure $K$ of $k$, such a norm-one frame
exists. Indeed $1^\perp$ is nondegenerate, so choose a nonisotropic pure
vector and rescale it to norm one over $K$. Its orthogonal complement with
$1$ is nondegenerate, allowing a similarly normalized $j$. The quaternion
frame fact makes $H$ nondegenerate; its four-dimensional orthogonal
complement supplies $\ell$ in the same way. The square roots are used only
in $K$, not assumed to lie in $k$. Over a finite field of odd order, the
retained binary-form argument also supplies these norm-one choices over the
field itself. Either route suffices for the nondegeneracy conclusion.

\section{A complete symbolic contraction calculation}

Take a norm-one frame and write
\[
 (e_1,e_2,e_3,e_4,e_5,e_6,e_7)
     =(i,j,ij,\ell,i\ell,j\ell,(ij)\ell).
\]
Its norm half-polar matrix on $V$ is the identity. Let $e^1,\ldots,e^7$
be the dual basis, write $e^{abc}=e^a\wedge e^b\wedge e^c$, and put
$\Omega=e^{1234567}$. The multiplication rule~\eqref{eq:double} gives
\begin{equation}\label{eq:form}
 \Phi=e^{123}+e^{145}-e^{167}+e^{246}+e^{257}+e^{347}-e^{356}.
\end{equation}
For example $i(j\ell)=(ji)\ell=-(ij)\ell$, giving the coefficient
$-1$ at $167$, while $j(i\ell)=(ij)\ell$ gives $+1$ at $257$.
The other five products are $ij=e_3$, $i\ell=e_5$, $j\ell=e_6$,
$(ij)\ell=e_7$ and $(ij)(i\ell)=-j\ell$. All other increasing triples
have zero coefficient, since the product of two distinct frame vectors is
the signed frame vector on their Cayley line. Thus~\eqref{eq:form}
specifies all35 increasing coefficients, with no numerical model input.

For any alternating three-form, set
\[
 b_\Phi(x,y)=(\iota_x\Phi)\wedge(\iota_y\Phi)\wedge\Phi
       \ \in\Lambda^7V^*.
\]
The following table gives every diagonal calculation. In each row the three
displayed terms of $\iota_{e_r}\Phi$ are called $A_r,B_r,C_r$ in that order.
The last three columns give the coefficient of $\Omega$ in
$A_r\wedge B_r\wedge\Phi$, $A_r\wedge C_r\wedge\Phi$ and
$B_r\wedge C_r\wedge\Phi$, respectively.
\[
\begin{array}{c|c|rrr}
 r&\iota_{e_r}\Phi&AB\Phi&AC\Phi&BC\Phi\\\hline
 1&e^{23}+e^{45}-e^{67}&-1&-1&-1\\
 2&-e^{13}+e^{46}+e^{57}&-1&-1&-1\\
 3&e^{12}+e^{47}-e^{56}&-1&-1&-1\\
 4&-e^{15}-e^{26}-e^{37}&-1&-1&-1\\
 5&e^{14}-e^{27}+e^{36}&-1&-1&-1\\
 6&e^{17}+e^{24}-e^{35}&-1&-1&-1\\
 7&-e^{16}+e^{25}+e^{34}&-1&-1&-1
\end{array}
\]
These entries involve only reordering seven distinct covectors. For the
first row the three contributions are
\[
 e^{2345}\wedge(-e^{167})=-\Omega,\qquad
 (-e^{2367})\wedge e^{145}=-\Omega,\qquad
 (-e^{4567})\wedge e^{123}=-\Omega.
\]
For any row, each of $A_r,B_r,C_r$ is decomposable and has square zero.
Two-forms commute in the exterior algebra. Hence the row gives
\[
 b_\Phi(e_r,e_r)
   =2(A_rB_r\Phi+A_rC_r\Phi+B_rC_r\Phi)=-6\Omega.
\]

For completeness the off-diagonal entries can be eliminated without a
further table or an appeal to the involution theorem being proved useful
here. Label the seven vectors by the nonzero degrees of $\mathbb F_2^3$:
\[
 (d_1,\ldots,d_7)=(100,010,110,001,101,011,111).
\]
Every triple in~\eqref{eq:form} has degree sum zero. For a linear
functional $\chi:\mathbb F_2^3\to\mathbb F_2$, the diagonal map
$D_\chi e_r=(-1)^{\chi(d_r)}e_r$ consequently preserves $\Phi$.
If $\chi\ne0$, exactly four of the seven degrees have value one, so
$\det D_\chi=(-1)^4=1$; the trivial functional gives determinant one
as well. Naturality of contraction and exterior product gives
\[
 b_\Phi(D_\chi x,D_\chi y)=(D_\chi^{-1})^*b_\Phi(x,y)=b_\Phi(x,y).
\]
This covariance identity uses no nondegeneracy assumption. If $r\ne s$,
choose $\chi$ with $\chi(d_r+d_s)=1$. Bilinearity then makes the left
side $-b_\Phi(e_r,e_s)$. As the characteristic is not two, that entry is
zero. We have proved the full identity
\begin{equation}\label{eq:metric}
                    b_\Phi=-6B\,\Omega
\end{equation}
in this exact norm-one frame, including its orientation and signs. Choosing
a different Cayley-frame convention can change the displayed volume sign;
it does not change the nondegeneracy conclusion.

\begin{theorem}\label{thm:nondegenerate}
Let $\C$ and $\Phi$ be as above. If $\operatorname{char}k\notin\{2,3\}$,
the tensor-valued contraction form $b_\Phi$ is nondegenerate. This holds
in particular for the normalized split Cayley cross-product construction
over $\mathbb F_{11}$.
\end{theorem}
\begin{proof}
Over the extension $K$, equation~\eqref{eq:metric} identifies the
contraction form with a nonzero scalar times the nonsingular norm form,
valued in a nonzero volume form. Thus its radical over $K$ is zero.
Contraction, exterior product and the associated linear map commute with
extension of scalars. Equivalently, the rank of any representing matrix is
unchanged by a field extension. A nonzero vector in the radical over $k$
would therefore remain a nonzero radical vector over $K$, a contradiction.
\end{proof}

The characteristic restriction is substantive. The same symbolic calculation
in characteristic three gives $b_\Phi=0$ for this normalized form, despite
nondegeneracy of its norm. Thus norm nondegeneracy alone was not silently
used as contraction nondegeneracy. The theorem supplies the missing
mathematical implication with its precise characteristic hypothesis.

\section{The retained normalization and the reviewed involution theorem}

In the retained criterion
\emph{A finite tensor certificate for Cayley normalization}, let $\beta$ be
the nondegenerate symmetric form, $\tau$ the alternating form, and $T$ its
$\beta$-dual cross product. A nonzero $\kappa$ with the complete double-cross
identity constructs
\[
 (a,x)(b,y)=(ab-\kappa\beta(x,y),ay+bx+T(x,y)),\qquad
 N(a,x)=a^2+\kappa\beta(x,x).
\]
Its pure norm half-polar form is $B=\kappa\beta$, and its normalized
three-form is exactly
\[
        B(T(x,y),z)=\kappa\tau(x,y,z).
\]
Consequently Theorem~\ref{thm:nondegenerate} applies to that normalized
form over $\mathbb F_{11}$. The previously accepted certificate has
$\kappa=7$; neither that value nor any actual coefficient is computed
here. If a three-form is rescaled by $c\ne0$, its contraction form is
rescaled by $c^3$, so its nondegeneracy is unchanged. A rescaling or change
of coordinates nevertheless needs an explicit exact subject join before
applying the assertion to a saved form.

\begin{corollary}\label{cor:projector}
For this normalized Cayley form over a field of characteristic different
from two and three, an idempotent $E$ of rank one gives an involution
$s=2E-I$ which does not preserve $\Phi$. If a subgroup $D$ preserves
$\Phi$, neither $Ds$ nor $sD$ meets its stabilizer.
\end{corollary}
\begin{proof}
The decomposition $V=\operatorname{im}E\oplus\ker E$ shows that $s$ is
a nonidentity involution with positive and negative eigenspace dimensions
one and six. Theorem~\ref{thm:nondegenerate} supplies the hypothesis of
the already reviewed note \emph{Involutions preserving an alternating form
in dimension seven}, exact7170-byte edition with SHA256
\texttt{8e16a4c31a114d8bbf0389b884a5f3634dd1843babbf310c52c7d0e9ea78b1aa}.
That theorem forces the dimensions three and four for a preserving
nonidentity involution, a contradiction. Its imported radical lemma and
proof are reused, not repeated. The stabilizer is a subgroup, so $ds$ or
$sd$ in it with $d\in D$ would put $s$ in it as well.
\end{proof}

\paragraph{Attribution and application boundary.}
The Cayley-frame and normalization constructions are borrowed from the
retained programme and its stated classical inputs. The contraction formula
is a standard relation for a Cayley three-form; this note supplies its full
symbolic calculation, conventions and conditional join, without a novelty
or priority claim. The reviewed involution argument arose in the Root
mathematical discussion and was written by /root/mathematical\_review.
This connecting proof was written by /root/forest\_resume and requires its
own distinct review.

For position2739, the previous Source already constructs its rank-one
idempotent and the exact $2E-I$, and directly certifies preservation of its
saved alternating form by the four marks. This note does not decode the
saved form, establish a new equality with the normalization list, restore
an unavailable projector/review body, or adopt a model. Its application
requires the separate exact accepted model/constructor/equality premise
discussed in \texttt{projector\_form\_scope.rst}. No new numerical role,
case, job, conjecture answer, manuscript, Lean proof or complexity rate
is claimed, and no accepted numerical certificate or registration Source
is changed.
\end{document}
